\lstset{
  language=Python,
  basicstyle=\ttfamily\footnotesize,
  numbers=left,
  numberstyle=\tiny,
  stepnumber=1,
  numbersep=5pt,
  backgroundcolor=\color{white},
  showspaces=false,
  showstringspaces=false,
  showtabs=false,
  frame=single,
  tabsize=2,
  captionpos=b,
  breaklines=true,
  breakatwhitespace=false,
  keywordstyle=\color{blue}\bfseries,
  commentstyle=\color{green!50!black}\textit,
  stringstyle=\color{red!70!black}
}

\appendix
\chapter{Source Code} \label{appendice:code}
This section contains the source code used to carry out the experiments described in the thesis. The files are divided by functionality.

All the complete files are also available in the repository attached to the project \cite{molinario2025thesis}.

\section{Dataset 2: Subsampling from Deep Firearm}
\begin{lstlisting}[caption=Python code for Subsampling from "DeepFirearm" dataset, label=lst:python_DB2]
import os
import shutil
import random

# === Source and destination paths ===
# Define the source directory containing the original dataset
root_dir = "/media/lello/Onedrive/20_Tesi/Dataset/DATASET SCELTI/02_httpsgithub.comjdhaodeep_firearm_COMPLETO/firearm-dataset-002/firearm-dataset/train"

# Define the destination directory where the selected subset will be saved
dest_dir = "/media/lello/Onedrive/20_Tesi/Dataset/DATASET SCELTI/02_httpsgithub.comjdhaodeep_firearm_COMPLETO/SubsetScelto/"

# Maximum number of total images to select across all classes
max_total_images = 1000

# Maximum number of images allowed per individual class
max_per_class = 200

# === Clean and prepare the destination directory ===
# If the destination folder already exists, it is removed and recreated to avoid duplication
if os.path.exists(dest_dir):
    shutil.rmtree(dest_dir)
os.makedirs(dest_dir, exist_ok=True)

# === List and sort all available class directories ===
# Extract all subdirectories representing class names in the source dataset
all_classes = sorted([d for d in os.listdir(root_dir) if os.path.isdir(os.path.join(root_dir, d))])
print(f" {len(all_classes)} available classes detected.")

# === Class selection and image sampling ===
total_images = 0  # Counter for total images added to the subset
selected_classes = []  # Track selected classes and their image counts

# Iterate over each class directory
for class_name in all_classes:
    src_dir = os.path.join(root_dir, class_name)

    # Collect all image files with valid extensions
    images = [f for f in os.listdir(src_dir) if f.lower().endswith(('.jpg', '.jpeg', '.png'))]

    # Skip classes with insufficient image count
    if len(images) < 20:
        continue

    # Select up to max_per_class images randomly
    n_select = min(len(images), max_per_class)
    sampled_images = random.sample(images, n_select)

    # Create a corresponding subfolder in the destination directory
    dest_class_dir = os.path.join(dest_dir, class_name)
    os.makedirs(dest_class_dir, exist_ok=True)

    # Copy the sampled images to the new destination
    for img in sampled_images:
        shutil.copy(os.path.join(src_dir, img), os.path.join(dest_class_dir, img))

    # Update counters and log the class inclusion
    total_images += len(sampled_images)
    selected_classes.append((class_name, len(sampled_images)))

    print(f" {class_name}: {len(sampled_images)} images selected")

    # Stop if the overall image quota has been reached
    if total_images >= max_total_images:
        break

# === Summary of the dataset subset ===
print(f"\n Total images selected: {total_images} across {len(selected_classes)} classes")
for cname, count in selected_classes:
    print(f" - {cname}: {count} images")

print(f"\n Final dataset created at: {dest_dir}")

}
\end{lstlisting}


\section{Dataset 3: Scraping Google Images}
\begin{lstlisting}[caption=Python Code Custom Collection via "Google Images" Scraping, label=lst:python_DB3]
from icrawler.builtin import GoogleImageCrawler
from PIL import Image
import os
import shutil
import random
import zipfile

# === CONFIGURATION PARAMETERS ===
max_total_images = 1000        # Upper limit for the total number of images to be collected
max_per_class = 125            # Maximum number of images allowed per individual category
output_base_dir = "dataset_scraping_google"
os.makedirs(output_base_dir, exist_ok=True)

# === CATEGORY DEFINITIONS (15 semantically realistic firearm-related classes) ===
categories = {
    "concealed_weapon": "concealed weapon",
    "gun_on_table": "gun on table",
    "person_with_gun": "person with gun",
    "gun_in_bag": "gun in bag",
    "toy_gun": "toy gun",
    "airsoft_weapon": "airsoft gun realistic",
    "holstered_gun": "holstered pistol",
    "damaged_weapon": "damaged firearm",
    "gun_in_drawer": "gun inside drawer",
    "gun_on_floor": "gun on floor",
    "person_pointing_gun": "person pointing gun",
    "gun_next_to_object": "gun next to phone",
    "gun_in_backpack": "gun inside backpack",
    "3d_printed_gun": "3d printed firearm",
    "gun_in_car": "gun in car seat",
}

# === FUNCTION: Download images from Google Images for a given category ===
def download_class_images(label, query, n_images):
    output_dir = os.path.join(output_base_dir, label)
    if os.path.exists(output_dir):
        print(f" Skipping '{label}', already exists.")
        return 0
    crawler = GoogleImageCrawler(storage={'root_dir': output_dir})
    crawler.crawl(keyword=query, max_num=n_images, min_size=(200, 200))  # Discard thumbnails
    return len(os.listdir(output_dir))

# === FUNCTION: Remove corrupted or unreadable image files ===
def remove_invalid_images():
    total_removed = 0
    for class_dir in os.listdir(output_base_dir):
        folder = os.path.join(output_base_dir, class_dir)
        for f in os.listdir(folder):
            try:
                Image.open(os.path.join(folder, f)).verify()
            except:
                os.remove(os.path.join(folder, f))
                total_removed += 1
    print(f" Removed {total_removed} invalid or corrupted images")

# === IMAGE COLLECTION PHASE ===
downloaded_total = 0
class_iter = list(categories.items())

# Iterate through the predefined categories until the total quota is met
while downloaded_total < max_total_images and class_iter:
    label, query = class_iter.pop(0)
    print(f"Downloading '{label}'...")
    n = download_class_images(label, query, max_per_class)
    downloaded_total += n
    print(f" {n} images downloaded for '{label}' (Running total: {downloaded_total})")

# === CLEANING PHASE ===
remove_invalid_images()

# === ZIP ARCHIVING PHASE ===
zip_path = "dataset_scraping_google.zip"
with zipfile.ZipFile(zip_path, 'w') as zipf:
    for root, _, files in os.walk(output_base_dir):
        for file in files:
            zipf.write(os.path.join(root, file),
                       os.path.relpath(os.path.join(root, file), output_base_dir))

print(f"\n Collection complete: {downloaded_total} images downloaded")
print(f" ZIP archive saved to: {zip_path}")


}


\end{lstlisting}



\section{Dataset 4: Scraping from YouTube}
\begin{lstlisting}[caption=Python code OSINT collection from YouTube, label=lst:python_DB4]
# === REQUIREMENTS ===
# This script collects images from Telegram channels,
# using Google search via SerpAPI to automatically identify relevant public channels.

import os
from telethon.sync import TelegramClient
from telethon.tl.types import MessageMediaPhoto
from telethon.errors import UsernameNotOccupiedError
from serpapi import GoogleSearch

# === TELEGRAM API PARAMETERS ===
# Replace with your own credentials obtained from https://my.telegram.org
api_id = "INSERT-YOUR-ID"
api_hash = 'INSERT-YOUR-KEY'
session_name = "anon"  # Saved session name (reused across runs)

# === SERPAPI PARAMETERS ===
# Replace with your valid SerpAPI key
serpapi_api_key = "INSERT-YOUR-KEY"

# === OUTPUT FOLDER FOR DOWNLOADED IMAGES ===
output_dir = "osint_telegram"
os.makedirs(output_dir, exist_ok=True)

# === FUNCTION: Search for Telegram channel names using SerpAPI and Google ===
def search_telegram_channels_serpapi(query, max_results=10):
    params = {
        "engine": "google",
        "q": f"site:t.me {query}",  # Restrict search to Telegram links
        "api_key": serpapi_api_key
    }

    search = GoogleSearch(params)
    results = search.get_dict()
    links = []

    if "organic_results" in results:
        for result in results["organic_results"]:
            link = result.get("link", "")
            if "t.me/" in link:
                channel = link.split("t.me/")[1].split("/")[0].split("?")[0]
                if channel not in links:
                    links.append(channel)
            if len(links) >= max_results:
                break
    return links

# === AUTOMATED CHANNEL DISCOVERY ===
search_terms = [
    "military footage",
    "firearm",
    "combat footage",
    "weapons and equipment archives"
]

channels = []
for term in search_terms:
    found = search_telegram_channels_serpapi(term, max_results=5)
    channels.extend(found)

# === MANUAL CHANNEL WHITELIST (VERIFIED SOURCES) ===
manual_channels = [
    "weaponsandequipmentarchives",
    "SmithandWessoninc47",
    "UAWeapons"
]

channels.extend(manual_channels)

# === FORMAT AND CLEAN CHANNEL LIST ===
channels = list(set([f"@{ch}" for ch in channels]))
print(f"\U0001F4E1 Channels discovered: {channels}")

# === TELEGRAM ACCESS AND MEDIA EXTRACTION ===
client = TelegramClient(session_name, api_id, api_hash)

with client:
    count_total = 0

    for channel in channels:
        count = 0
        channel_name = channel.replace("@", "")
        channel_dir = os.path.join(output_dir, channel_name)
        os.makedirs(channel_dir, exist_ok=True)
        print(f"\nProcessing channel: {channel}")

        try:
            for msg in client.iter_messages(channel, limit=200):
                if isinstance(msg.media, MessageMediaPhoto):
                    filename = f"telegram_{count}.jpg"
                    filepath = os.path.join(channel_dir, filename)
                    # Avoid re-downloading existing images
                    if not os.path.exists(filepath):
                        client.download_media(msg.media, file=filepath)
                        count += 1
                    else:
                        count += 1
            print(f" {channel_name}: {count} images saved or already present")
            count_total += count

        except UsernameNotOccupiedError:
            print(f"Error in channel '{channel_name}': username does not exist")

        except Exception as e:
            print(f"Error in channel '{channel_name}': {e}")

    print(f"\n--- Extraction complete ---")
    print(f"\n Total Telegram images saved: {count_total}")

}

\end{lstlisting}






\section{Dataset 4: Scraping OSINT from Telegram}
\begin{lstlisting}[caption=Python Code OSINT Collection from Telegram, label=lst:python_DB4]
# === REQUIREMENTS ===
# This script collects visual samples from publicly accessible YouTube videos and,
# in case of failure, falls back to Google Images.
# It is intended for building realistic OSINT-based datasets with weapon-related visual content.

from icrawler.builtin import GoogleImageCrawler
import os
import urllib.request
from PIL import Image
import shutil
import yt_dlp

# === BASE CONFIGURATION ===

# Output directory where all downloaded images will be stored
output_base_dir = "osint_youtube"

# Maximum number of videos to retrieve per search query
max_videos = 100

# Categories (semantic labels) mapped to search queries
# Each query corresponds to a realistic forensic or OSINT context
categories = {
    "cctv_firearm_incident": "firearm incident CCTV",
    "gun_in_bag": "gun in bag",
    "toy_gun_or_airsoft": "airsoft toy gun",
    "training_police_shooting": "police shooting training",
    "person_with_gun": "person with gun",
    "crime_scene_weapon": "crime scene weapon"
}

# === FALLBACK: Google Images ===
# If no thumbnail is retrieved from YouTube, use Google Image Search as backup
def fallback_google_images(query, output_dir, max_images=100):
    crawler = GoogleImageCrawler(storage={'root_dir': output_dir})
    crawler.crawl(keyword=query, max_num=max_images, min_size=(200, 200))
    print(f" Google fallback completed: {max_images} images saved in {output_dir}")

# === FUNCTION: Download video thumbnails using yt-dlp ===
def download_youtube_thumbnails(query, label):
    output_dir = os.path.join(output_base_dir, label)

    # Remove existing folder to avoid duplication or overwriting
    if os.path.exists(output_dir):
        shutil.rmtree(output_dir)
    os.makedirs(output_dir, exist_ok=True)

    ydl_opts = {
        'quiet': True,             # Suppress console output
        'skip_download': True,     # Do not download the actual videos
        'extract_flat': False      # Allow access to metadata and thumbnails
    }

    count = 0

    try:
        with yt_dlp.YoutubeDL(ydl_opts) as ydl:
            result = ydl.extract_info(f"ytsearch{max_videos}:{query}", download=False)
            entries = result.get('entries', [])
            for i, entry in enumerate(entries):
                thumb_url = entry.get('thumbnail')
                if thumb_url:
                    try:
                        filename = os.path.join(output_dir, f"{label}_{i}.jpg")
                        urllib.request.urlretrieve(thumb_url, filename)

                        # Validate image integrity
                        Image.open(filename).verify()
                        count += 1
                    except:
                        # Remove corrupt or unreadable file
                        if os.path.exists(filename):
                            os.remove(filename)

    except Exception as e:
        print(f"yt-dlp error: {e}")

    # If no thumbnails were retrieved, fall back to Google
    if count == 0:
        print(f" No images found for '{label}', switching to Google fallback...")
        fallback_google_images(query, output_dir, max_images=max_videos)

    return count

# === MAIN SCRIPT EXECUTION ===
# Clear previous output directory and create a fresh one
if os.path.exists(output_base_dir):
    shutil.rmtree(output_base_dir)
os.makedirs(output_base_dir, exist_ok=True)

# Iterate through all categories and download relevant images
total_downloaded = 0
for label, query in categories.items():
    print(f"\nDownloading: {query}")
    n = download_youtube_thumbnails(query, label)
    print(f" {label}: {n} images saved")
    total_downloaded += n

print(f"\n Total images saved: {total_downloaded}")


}

\end{lstlisting}






\section{Dataset 5: Scraping Deep Web}
\begin{lstlisting}[caption=Python Code Collected from the Deep Web, label=lst:python_DB5]
# === DEEP WEB IMAGE EXTRACTION PIPELINE ===
# This script performs automated crawling of .onion websites indexed by Ahmia,
# leveraging the Tor network to access hidden services and retrieve weapon-related imagery.
# It includes quality control, deduplication, and filtering heuristics.

import os
import requests
from bs4 import BeautifulSoup
from urllib.parse import urljoin, unquote, urlparse
import time
from PIL import Image
from io import BytesIO
import hashlib

# === PROXY CONFIGURATION FOR TOR ===
proxies = {
    'http': 'socks5h://127.0.0.1:9050',   # Ensure Tor daemon is running on port 9050
    'https': 'socks5h://127.0.0.1:9050'
}

# === OUTPUT DIRECTORY SETUP ===
output_dir = "deepweb_dataset"
os.makedirs(output_dir, exist_ok=True)

# === QUERY LIST ===
# A set of semantically relevant keywords for weapon-related content
search_queries = [
    "gun marketplace", "firearm photo", "pistol images", "smuggled weapons",
    "AK-47 picture", "gun trading", "assault rifle image", "sniper darkweb",
    "hidden firearm photo", "weapon", "gun", "pistol", "AK-47", "rifle", "sniper", "firearm"
]

# === FUNCTION: Extract .onion URLs from Ahmia ===
def search_ahmia(query):
    search_url = f"https://ahmia.fi/search/?q={query.replace(' ', '+')}"
    try:
        res = requests.get(search_url, timeout=15)
        soup = BeautifulSoup(res.text, "html.parser")
        links = []
        for a in soup.find_all("a", href=True):
            href = a['href']
            if "/redirect?search_term=" in href:
                redirect_url = unquote(href.split("redirect_url=")[-1])
                if ".onion" in redirect_url and redirect_url not in links:
                    links.append(redirect_url)
        links.sort(key=lambda l: len(l), reverse=True)  # Prioritize longer URLs
        print(f" Found {len(links)} .onion links for '{query}'")
        return links
    except Exception as e:
        print(f"Ahmia search error for '{query}': {e}")
        return []

# === FUNCTION: Download and validate images from a given .onion page ===
def download_images_from_onion(url, session, output_subdir, max_images=15):
    try:
        res = session.get(url, timeout=15)
        soup = BeautifulSoup(res.text, "html.parser")
        imgs = soup.find_all("img")
        count = 0

        # Skip common noise elements
        ban_keywords = ["logo", "banner", "avatar", "icon", "sprite", "captcha"]

        for img in imgs:
            if count >= max_images:
                break
            src = img.get("src")
            if not src or any(keyword in src.lower() for keyword in ban_keywords):
                continue
            img_url = urljoin(url, src)

            try:
                r = session.get(img_url, timeout=10)
                if r.status_code == 200 and "image" in r.headers.get("Content-Type", ""):
                    content_hash = hashlib.md5(r.content).hexdigest()
                    ext = r.headers.get("Content-Type").split("/")[-1].split(";")[0]
                    filename = f"{content_hash}.{ext}"
                    filepath = os.path.join(output_subdir, filename)

                    # Avoid saving duplicates or small artifacts
                    if not os.path.exists(filepath):
                        img_obj = Image.open(BytesIO(r.content))
                        if img_obj.width < 200 or img_obj.height < 200:
                            continue
                        with open(filepath, "wb") as f:
                            f.write(r.content)
                        count += 1
            except Exception:
                continue
        return count
    except Exception as e:
        print(f"Error visiting .onion site '{url}': {e}")
        return 0

# === INITIALIZE TOR SESSION ===
session = requests.Session()
session.proxies.update(proxies)

# === MAIN LOOP: Iterate over queries and sites ===
total_downloaded = 0
visited_paths = set()  # Used to skip duplicate paths with different domain prefixes

for query in search_queries:
    print(f"\n Query: {query}")
    sub_dir = os.path.join(output_dir, query.replace(" ", "_"))
    os.makedirs(sub_dir, exist_ok=True)

    onion_links = search_ahmia(query)

    for onion_url in onion_links:
        parsed = urlparse(onion_url)
        url_path_key = f"{parsed.netloc}{parsed.path}"

        if url_path_key in visited_paths:
            print(f" Skipping duplicate path: {onion_url}")
            continue

        visited_paths.add(url_path_key)
        print(f" Visiting: {onion_url}")
        n = download_images_from_onion(onion_url, session, sub_dir)
        print(f" Images saved: {n}")
        total_downloaded += n
        time.sleep(3)  # polite delay between requests

print(f"\n Deep web scraping completed. Total images saved: {total_downloaded}")

\end{lstlisting}





\section{Dataset: Image Selection for Final Subset}
\begin{lstlisting}[caption={Python script for selecting images for robustness analysis}, label={lst:python_DBSubSetChoice}]

    import os
    import csv
    import random
    import shutil
    from PIL import Image
    
    # Accepted file extensions for image files
    VALID_EXTENSIONS = {'.jpg', '.jpeg', '.png', '.bmp', '.gif', '.webp'}
    
    def is_valid_image(filepath, min_size=300):
        """
        Checks whether the file is a valid image with at least the minimum required dimensions.
        """
        try:
            with Image.open(filepath) as img:
                width, height = img.size
                if width >= min_size and height >= min_size:
                    return width, height
        except Exception:
            pass
        return None
    
    def normalize_dataset_name(name):
        """
        Cleans the sub-dataset name by removing descriptions in parentheses or URLs.
        """
        return name.replace(" (httpswww.kaggle.comdatasetssnehilsanyalweapon-detection-test)", "")\
                   .replace("httpsgithub.comjdhaodeep_firearm_COMPLETO", "")\
                   .strip()
    
    def select_and_copy_images(dataset_paths, output_dir, min_size=300, max_images=200):
        """
        Extracts valid images from the specified sub-datasets, copies them into output_dir,
        and writes a summary CSV file. Skips datasets already complete.
    
        Args:
            dataset_paths (dict): Mapping between logical names and folder paths.
            output_dir (str): Final output directory.
            min_size (int): Minimum image dimension (in pixels).
            max_images (int): Maximum number of images to select per dataset.
        """
        os.makedirs(output_dir, exist_ok=True)
        summary_path = os.path.join(output_dir, "selected_images_summary.csv")
        all_rows = []
    
        fieldnames = ["dataset", "original_path", "copied_path", "width", "height", "size_bytes", "is_real_weapon"]
    
        for dataset_label, dataset_root in dataset_paths.items():
            print(f"\nAnalyzing dataset: {dataset_label}")
    
            clean_name = normalize_dataset_name(dataset_label)
            dest_dir = os.path.join(output_dir, clean_name)
            os.makedirs(dest_dir, exist_ok=True)
    
            existing = [f for f in os.listdir(dest_dir) if os.path.splitext(f)[1].lower() in VALID_EXTENSIONS]
            if len(existing) >= max_images:
                print(f"  Found {len(existing)} images in '{dest_dir}', skipping.")
                continue
    
            valid_images = []
            for root, _, files in os.walk(dataset_root):
                for f in files:
                    ext = os.path.splitext(f)[1].lower()
                    if ext not in VALID_EXTENSIONS:
                        continue
                    full_path = os.path.join(root, f)
                    dim = is_valid_image(full_path, min_size)
                    if dim:
                        width, height = dim
                        valid_images.append((full_path, width, height))
    
            print(f" Found {len(valid_images)} valid images.")
            selected = random.sample(valid_images, min(len(valid_images), max_images))
    
            for i, (src, w, h) in enumerate(selected):
                ext = os.path.splitext(src)[1].lower()
                filename = f"{clean_name}_{i+1:03d}{ext}"
                dst = os.path.join(dest_dir, filename)
                shutil.copy2(src, dst)
    
                all_rows.append({
                    "dataset": clean_name,
                    "original_path": src,
                    "copied_path": dst,
                    "width": w,
                    "height": h,
                    "size_bytes": os.path.getsize(src),
                    "is_real_weapon": True
                })
    
        # Write summary CSV
        with open(summary_path, "w", newline='', encoding="utf-8") as f:
            writer = csv.DictWriter(f, fieldnames=fieldnames)
            writer.writeheader()
            writer.writerows(all_rows)
    
        print(f"\nTotal images copied: {len(all_rows)}")
        print(f" CSV saved to: {summary_path}")
    
    
    dataset_paths = {
        "01_archive": "../Dataset/01_archive (httpswww.kaggle.comdatasetssnehilsanyalweapon-detection-test)",
        "02_github_deep_firearm": "../Dataset/02_httpsgithub.comjdhaodeep_firearm_COMPLETO",
        "03_scraped_firearm_dataset": "../Dataset/03_scraped_firearm_dataset",
        "04_Osint_yt_telegram": "../Dataset/04_Osint_yt_telegram",
        "05_deepweb_dataset": "../Dataset/05_deepweb_dataset"
    }
    
    select_and_copy_images(dataset_paths, "../final_test_subset")
    

\end{lstlisting}



\section{Discovery of Out-Of-Distribution (OOD) Images}
\begin{lstlisting}[caption={Discovery of Out-Of-Distribution (OOD) Images}, label={appendix:clip-ood-script}]

    import os
    import torch
    import clip
    import pandas as pd
    from PIL import Image
    import matplotlib.pyplot as plt
    from tqdm import tqdm
    
    image_extensions = {'.jpg', '.jpeg', '.png', '.bmp', '.gif', '.webp'}
    marked_odd = set()
    current_images = []
    
    def extract_subdataset(path):
        """
        Extracts the sub-dataset name from the image path.
        """
        for part in os.path.normpath(path).split(os.sep):
            if part.startswith(("01_", "02_", "03_", "04_", "05_")):
                return part
        return "unknown"
    
    def run_clip_and_annotation_pipeline(image_dir, batch_size=50):
        """
        Automatically classifies images using CLIP, enables manual annotation, and saves final CSV results.
    
        Args:
            image_dir (str): Directory containing images organized by sub-dataset.
            batch_size (int): Number of images to display per screen.
        """
        # === CLIP: automatic classification
        print("\n Starting CLIP classification...")
        device = "cuda" if torch.cuda.is_available() else "cpu"
        model, preprocess = clip.load("ViT-B/32", device=device)
    
        prompts = [
            "a photo of a handgun",
            "a photo of a rifle",
            "a photo of a toy gun",
            "a photo of a person",
            "a photo of a household object"
        ]
    
        text_inputs = clip.tokenize(prompts).to(device)
        text_features = model.encode_text(text_inputs)
    
        clip_results = []
        all_images = []
    
        for root, _, files in os.walk(image_dir):
            for file in files:
                if os.path.splitext(file)[1].lower() in image_extensions:
                    path = os.path.join(root, file)
                    all_images.append(path)
                    try:
                        image = preprocess(Image.open(path)).unsqueeze(0).to(device)
                        image_features = model.encode_image(image)
                        similarity = (image_features @ text_features.T).softmax(dim=-1).detach().cpu().numpy()[0]
                        best = prompts[similarity.argmax()]
                        label = "OOD" if best not in ["a photo of a handgun", "a photo of a rifle"] else "in-distribution"
    
                        clip_results.append({
                            "image_path": path,
                            "clip_label": label,
                            "clip_match": best,
                            "clip_confidence": round(float(similarity.max()), 4),
                            "subdataset": extract_subdataset(path)
                        })
                    except Exception as e:
                        print(f" Error on {path}: {e}")
    
        clip_df = pd.DataFrame(clip_results)
        clip_csv = os.path.join(image_dir, "clip_ood_results.csv")
        clip_df.to_csv(clip_csv, index=False)
        print(f" CLIP results saved to: {clip_csv}")
    
        # === Manual annotation
        print("\nManual annotation...")
    
        clip_ood_set = set(clip_df[clip_df["clip_label"] == "OOD"]["image_path"])
    
        def onclick(event):
            if event.inaxes:
                idx = int(event.inaxes.get_label())
                path = current_images[idx]
                if path in marked_odd:
                    marked_odd.remove(path)
                else:
                    marked_odd.add(path)
                update_title(event.inaxes, path)
                event.canvas.draw_idle()
    
        def update_title(ax, path):
            is_weapon = path not in marked_odd
            is_ood = path in clip_ood_set
    
            label = "WEAPON" if is_weapon else "NON-WEAPON"
            if is_ood:
                label += " [CLIP: NON-WEAPON]"
            ax.set_title(label, fontsize=8)
    
            # Border color
            color = 'green' if is_weapon else 'red'
            for side in ax.spines:
                ax.spines[side].set_color(color)
                ax.spines[side].set_linewidth(2)
    
            # Remove previous icons
            if hasattr(ax, "_icon"):
                ax._icon.remove()
    
            # Add visual symbol
            symbol = "[OK]" if is_weapon else "[X]"
            ax._icon = ax.text(
                0.02, 0.85, symbol, transform=ax.transAxes,
                fontsize=20, color=color, verticalalignment='top', horizontalalignment='left',
                bbox=dict(facecolor='white', edgecolor=color, boxstyle='round,pad=0.3', alpha=0.7)
            )
    
        def show_thumbnails(batch, cols=6, size=3):
            global current_images
            current_images = batch
            rows = (len(batch) + cols - 1) // cols
            fig, axes = plt.subplots(rows, cols, figsize=(cols*size, rows*size))
            axes = axes.flatten()
    
            for i, path in enumerate(batch):
                img = Image.open(path)
                axes[i].imshow(img)
                axes[i].axis("off")
                axes[i].set_label(str(i))
                update_title(axes[i], path)
    
            for ax in axes[len(batch):]:
                ax.remove()
    
            fig.canvas.mpl_connect('button_press_event', onclick)
            plt.suptitle("Click to label as NON-WEAPON", fontsize=14)
            plt.tight_layout()
            plt.show()
    
        all_images.sort(key=lambda x: (x not in clip_ood_set, x))
    
        for i in range(0, len(all_images), batch_size):
            show_thumbnails(all_images[i:i+batch_size])
    
        manual_annotations = []
        for path in all_images:
            manual_annotations.append({
                "image_path": path,
                "manual_is_real_weapon": path not in marked_odd,
                "subdataset": extract_subdataset(path)
            })
    
        manual_df = pd.DataFrame(manual_annotations)
        manual_csv = os.path.join(image_dir, "manual_annotation_is_real_weapon.csv")
        manual_df.to_csv(manual_csv, index=False)
        print(f" Manual annotations saved to: {manual_csv}")
    
        # === Final merge and summary statistics
        print("\n Merging automatic and manual results...")
        merged_df = pd.merge(clip_df, manual_df, on=["image_path", "subdataset"])
        merged_csv = os.path.join(image_dir, "merged_clip_manual_results.csv")
        merged_df.to_csv(merged_csv, index=False)
    
        print(f" Merged file saved to: {merged_csv}")
    
        # Summary statistics
        n_auto_ood = (merged_df["clip_label"] == "OOD").sum()
        n_manual_ood = (~merged_df["manual_is_real_weapon"]).sum()
    
        summary_txt = os.path.join(image_dir, "ood_summary_stats.txt")
        with open(summary_txt, "w") as f:
            f.write("Summary of NON-WEAPON (OOD) images\n")
            f.write(f"Automatic (CLIP): {n_auto_ood}\n")
            f.write(f"Manual (user): {n_manual_ood}\n")
    
        print(f" Automatic CLIP: {n_auto_ood} NON-WEAPON images")
        print(f" Manual: {n_manual_ood} NON-WEAPON images")
        print(f" Statistics saved to: {summary_txt}")
    
    
    run_clip_and_annotation_pipeline("../final_test_subset")
    
\end{lstlisting}




\appendix
\chapter{Codice per la classificazione automatica con CLIP}
\label{appendix:clip-ood-script}

\section{Pipeline CLIP ViT-B/32 per la classificazione zero-shot}
\label{appendix:clip-section}

Questo script Python è stato sviluppato per effettuare la classificazione automatica di immagini tramite il modello \textbf{CLIP ViT-B/32}, in modalità \textit{zero-shot}. Il sistema analizza le immagini di input e le confronta con un set di prompt testuali predefiniti per identificare se esse rappresentano armi reali (\textit{in-distribution}) o oggetti fuori dominio (\textit{OOD}), come giocattoli, persone o oggetti di uso comune.

\subsection*{Funzionalità principali}
\begin{itemize}
    \item Caricamento del modello CLIP ViT-B/32 da OpenAI;
    \item Preprocessing delle immagini e tokenizzazione dei prompt testuali;
    \item Calcolo della similarità tra embedding visivi e testuali;
    \item Classificazione automatica in \texttt{in-distribution} o \texttt{OOD} sulla base del prompt più compatibile;
    \item Salvataggio dei risultati in formato CSV.
\end{itemize}

\noindent Il listato completo è riportato di seguito (Listato~\ref{lst:clip-script}).

\begin{lstlisting}[language=Python, caption={Script Python per la classificazione automatica con CLIP ViT-B/32}, label={lst:clip-script}, breaklines=true, basicstyle=\footnotesize\ttfamily]
import os
import torch
import clip
import pandas as pd
from PIL import Image

image_extensions = {'.jpg', '.jpeg', '.png', '.bmp', '.webp'}
result_dir = "../results"

def extract_subdataset(path):
    for part in os.path.normpath(path).split(os.sep):
        if part.startswith(("01_", "02_", "03_", "04_", "05_")):
            return part
    return "unknown"

def run_clip_classification(image_dir):
    print("\n Starting CLIP classification...")
    device = "cuda" if torch.cuda.is_available() else "cpu"
    model, preprocess = clip.load("ViT-B/32", device=device)

    prompts = [
        "a photo of a handgun",
        "a photo of a rifle",
        "a photo of a toy gun",
        "a photo of a person",
        "a photo of a household object"
    ]

    text_inputs = clip.tokenize(prompts).to(device)
    text_features = model.encode_text(text_inputs)

    clip_results = []

    for root, _, files in os.walk(image_dir):
        for file in files:
            if os.path.splitext(file)[1].lower() in image_extensions:
                path = os.path.join(root, file)
                try:
                    image = preprocess(Image.open(path)).unsqueeze(0).to(device)
                    image_features = model.encode_image(image)
                    similarity = (image_features @ text_features.T).softmax(dim=-1).detach().cpu().numpy()[0]
                    best = prompts[similarity.argmax()]
                    label = "OOD" if best not in ["a photo of a handgun", "a photo of a rifle"] else "in-distribution"

                    clip_results.append({
                        "image_path": path,
                        "clip_label": label,
                        "clip_match": best,
                        "clip_confidence": round(float(similarity.max()), 4),
                        "subdataset": extract_subdataset(path)
                    })
                except Exception as e:
                    print(f" Error on {path}: {e}")

    clip_df = pd.DataFrame(clip_results)
    clip_csv = os.path.join(result_dir, "clip_ood_results.csv")
    clip_df.to_csv(clip_csv, index=False)
    print(f" CLIP results saved to: {clip_csv}")

run_clip_classification("../Test_Dataset")
\end{lstlisting}



\chapter{Codice per la generazione di descrizioni e classificazione con BLIP}
\label{appendix:blip-ood-script}

\section{Pipeline BLIP per descrizione automatica e classificazione semantica}
\label{appendix:blip-section}

Il seguente script Python utilizza il modello \textbf{BLIP (Bootstrapped Language-Image Pretraining)} per generare in automatico descrizioni testuali delle immagini in input. Le caption ottenute vengono poi analizzate mediante un semplice sistema di filtraggio per determinare se l’immagine contiene un’arma reale (\texttt{in-distribution}) oppure è un contenuto fuori dominio (\texttt{OOD}).

\subsection*{Funzionalità principali}
\begin{itemize}
    \item Caricamento del modello BLIP da \texttt{Salesforce/blip-image-captioning-base};
    \item Generazione della descrizione testuale (caption) per ogni immagine;
    \item Classificazione semantica tramite ricerca di parole chiave nella caption;
    \item Salvataggio dei risultati in un file \texttt{CSV} con le colonne: percorso immagine, didascalia generata, etichetta, subdataset;
    \item Output riepilogativo con numero di immagini OOD e in-distribution.
\end{itemize}

\noindent Il codice completo è riportato nel Listato~\ref{lst:blip-script}.

\begin{lstlisting}[language=Python, caption={Script Python per la generazione di descrizioni e classificazione con BLIP}, label={lst:blip-script}, breaklines=true, basicstyle=\footnotesize\ttfamily]
import os
import pandas as pd
from PIL import Image
from tqdm import tqdm
import torch
from transformers import BlipProcessor, BlipForConditionalGeneration

class BLIPClassifier:
    def __init__(self, image_dir, result_dir="../results"):
        self.image_dir = image_dir
        self.result_dir = result_dir
        os.makedirs(self.result_dir, exist_ok=True)

        self.image_extensions = {'.jpg', '.jpeg', '.png', '.bmp', '.gif', '.webp'}
        self.device = torch.device("cuda" if torch.cuda.is_available() else "cpu")

        print("Caricamento modello BLIP...")
        self.processor = BlipProcessor.from_pretrained("Salesforce/blip-image-captioning-base")
        self.model = BlipForConditionalGeneration.from_pretrained(
            "Salesforce/blip-image-captioning-base"
        ).to(self.device)

        self.results = []

    def extract_subdataset(self, path):
        for part in os.path.normpath(path).split(os.sep):
            if part.startswith(("01_", "02_", "03_", "04_", "05_")):
                return part
        return "unknown"

    def classify_caption(self, caption):
        keywords = ["gun", "rifle", "weapon", "firearm", "pistol"]
        return "in-distribution" if any(k in caption.lower() for k in keywords) else "OOD"

    def run_inference(self):
        print(" Avvio inferenza immagini con BLIP...")
        for root, _, files in os.walk(self.image_dir):
            for fname in tqdm(files):
                if os.path.splitext(fname)[1].lower() in self.image_extensions:
                    img_path = os.path.join(root, fname)
                    try:
                        raw_image = Image.open(img_path).convert("RGB")
                        inputs = self.processor(raw_image, return_tensors="pt").to(self.device)
                        out = self.model.generate(**inputs)
                        caption = self.processor.decode(out[0], skip_special_tokens=True)

                        label = self.classify_caption(caption)

                        self.results.append({
                            "image_path": img_path,
                            "blip_caption": caption,
                            "blip_label": label,
                            "subdataset": self.extract_subdataset(img_path)
                        })
                    except Exception as e:
                        print(f"[!] Errore su {img_path}: {e}")

        self.save_results()

    def save_results(self):
        df = pd.DataFrame(self.results)
        output_csv = os.path.join(self.result_dir, "blip_ood_results.csv")
        df.to_csv(output_csv, index=False)
        print(f" File BLIP salvato in: {output_csv}")

        stats_txt = os.path.join(self.result_dir, "blip_summary_stats.txt")
        n_ood = (df["blip_label"] == "OOD").sum()
        n_indist = (df["blip_label"] == "in-distribution").sum()
        with open(stats_txt, "w") as f:
            f.write("BLIP Captioning Classification Summary\n")
            f.write(f"Total images: {len(df)}\n")
            f.write(f"Non-Weapon (OOD): {n_ood}\n")
            f.write(f"Weapon (in-distribution): {n_indist}\n")
        print(f" Statistics saved to: {stats_txt}")

# Esempio di utilizzo
if __name__ == "__main__":
    classifier = BLIPClassifier(image_dir="../Test_Dataset")
    classifier.run_inference()
\end{lstlisting}



\chapter{Codice per la classificazione binaria con ResNet18 + SVM}
\label{appendix:resnet18-svm-script}

\section{Pipeline ResNet18 + SVM per classificazione supervisata}
\label{appendix:resnet-section}

Questo script Python combina l’uso di una \textbf{ResNet18} pre-addestrata come estrattore di feature con un classificatore \textbf{SVM} (Support Vector Machine) per la classificazione binaria delle immagini in “arma” o “non arma”. L’architettura consente una separazione tra estrazione di feature e classificazione, utile in contesti forensi per maggiore interpretabilità.

\subsection*{Funzionalità principali}
\begin{itemize}
    \item Caricamento di una ResNet18 pre-addestrata da ImageNet;
    \item Estrazione degli embedding visivi per ciascuna immagine;
    \item Normalizzazione delle feature e classificazione tramite un classificatore SVM;
    \item Generazione di un file CSV con predizione binaria e probabilità associate;
    \item Classificazione in \texttt{in-distribution} o \texttt{OOD} in base alla soglia.
\end{itemize}

\noindent Il codice completo è riportato nel Listato~\ref{lst:resnet-svm-script}.

\begin{lstlisting}[language=Python, caption={Script Python per classificazione con ResNet18 + SVM}, label={lst:resnet-svm-script}, breaklines=true, basicstyle=\footnotesize\ttfamily]
import os
import torch
import pandas as pd
import numpy as np
from PIL import Image
from tqdm import tqdm
from sklearn.svm import SVC
from sklearn.preprocessing import StandardScaler
from torchvision import models, transforms

# === CONFIG ===
image_dir = "../Test_Dataset"
output_csv = os.path.join(image_dir, "resnet_ood_results.csv")
image_extensions = {'.jpg', '.jpeg', '.png', '.bmp', '.gif', '.webp'}
device = torch.device("cuda" if torch.cuda.is_available() else "cpu")

# === MODELLO ===
resnet = models.resnet18(weights=None)
resnet.load_state_dict(torch.load(r"/media/lello/Dati/Thesis/Attack/models/resnet18-f37072fd.pth", map_location="cpu"))
resnet = torch.nn.Sequential(*list(resnet.children())[:-1])  # Rimuove l'ultimo layer FC
resnet.eval().to(device)

# === FEATURE EXTRACTION TRANSFORM ===
preprocess = transforms.Compose([
    transforms.Resize((224, 224)),
    transforms.ToTensor(),
    transforms.Normalize(mean=[0.485, 0.456, 0.406],
                         std=[0.229, 0.224, 0.225]),
])

# === DUMMY TRAINING SVM SU EMBEDDING ARTIFICIALI ===
#  IN PRODUZIONE, ADDESTRA SU DATI REALI
X_train = np.random.randn(100, 512)
y_train = np.array([1]*50 + [0]*50)
scaler = StandardScaler().fit(X_train)
clf = SVC(probability=True).fit(scaler.transform(X_train), y_train)

# === INFERENZA ===
results = []

def extract_subdataset(path):
    for part in os.path.normpath(path).split(os.sep):
        if part.startswith(("01_", "02_", "03_", "04_", "05_")):
            return part
    return "unknown"

print("Estrazione feature ResNet + classificazione SVM...")
for root, _, files in os.walk(image_dir):
    for fname in tqdm(files):
        if os.path.splitext(fname)[1].lower() in image_extensions:
            path = os.path.join(root, fname)
            try:
                image = preprocess(Image.open(path).convert("RGB")).unsqueeze(0).to(device)
                with torch.no_grad():
                    feat = resnet(image).squeeze().cpu().numpy().flatten()

                prob = clf.predict_proba(scaler.transform([feat]))[0][1]
                label = "in-distribution" if prob > 0.5 else "OOD"

                results.append({
                    "image_path": path,
                    "resnet_label": label,
                    "resnet_confidence": round(prob, 4),
                    "subdataset": extract_subdataset(path)
                })
            except Exception as e:
                print(f"[!] Errore su {path}: {e}")

df = pd.DataFrame(results)
df.to_csv(output_csv, index=False)
print(f" File salvato in: {output_csv}")
\end{lstlisting}


\chapter{Codice per classificazione binaria con EfficientNet-B0}
\label{appendix:efficientnet-script}

\section{Training e inferenza con EfficientNet-B0}
\label{appendix:efficientnet-section}

In questo lavoro è stato impiegato il modello \textbf{EfficientNet-B0}, scelto per la sua efficienza computazionale e capacità di generalizzazione. Il modello è stato addestrato su un dataset bilanciato (arma / non arma), ottenuto da immagini reali suddivise automaticamente in cartelle \texttt{train/} e \texttt{val/}, come descritto nel file di preprocessing.

\subsection*{Fasi del processo}
\begin{enumerate}
  \item \textbf{Preparazione del dataset}: le immagini sono state organizzate per classe in cartelle separate, quindi divise automaticamente in training e validation set (80/20) tramite script Python (\texttt{04\_prepare\_dataset\_CNN.py});
  \item \textbf{Addestramento}: il modello EfficientNet-B0 è stato caricato con pesi pre-addestrati su ImageNet, e fine-tuned per 10 epoche su immagini arma / non arma (script \texttt{04\_efficientnet\_train\_CNN.py});
  \item \textbf{Inferenza}: il modello addestrato è stato utilizzato per classificare nuove immagini e determinare se sono \texttt{in-distribution} (arma) o \texttt{OOD} (non arma), confrontando le classi predette con una lista di keyword semantiche (script \texttt{04\_efficientnet\_ood\_test\_CNN.py}).
\end{enumerate}

\noindent I file completi sono riportati di seguito:

\subsection*{Preparazione del dataset}
\begin{lstlisting}[language=Python, caption={Script per suddivisione automatica del dataset in train/val}, label={lst:efficientnet-dataset}, breaklines=true, basicstyle=\footnotesize\ttfamily]
# Estratto da: 04_prepare_dataset_CNN.py
import os
import random
import shutil
from collections import defaultdict

source_dir = r"/Dataset_completo"
dest_dir = r"D:/Thesis/AI Tool Evaluation/efficientnet_dataset"
split_ratio = 0.8

for split in ["train", "val"]:
    os.makedirs(os.path.join(dest_dir, split), exist_ok=True)

class_files = defaultdict(list)
for root, _, files in os.walk(source_dir):
    for fname in files:
        if fname.lower().endswith((".jpg", ".jpeg", ".png", ".bmp", ".gif", ".webp")):
            full_path = os.path.join(root, fname)
            label = os.path.basename(root).lower()
            class_files[label].append(full_path)

for label, files in class_files.items():
    random.shuffle(files)
    split_point = int(len(files) * split_ratio)
    train_files = files[:split_point]
    val_files = files[split_point:]

    for split, split_files in [("train", train_files), ("val", val_files)]:
        class_dir = os.path.join(dest_dir, split, label)
        os.makedirs(class_dir, exist_ok=True)
        for f in split_files:
            shutil.copy2(f, os.path.join(class_dir, os.path.basename(f)))
\end{lstlisting}

\subsection*{Addestramento del modello EfficientNet-B0}
\begin{lstlisting}[language=Python, caption={Script per addestramento supervisionato di EfficientNet-B0}, label={lst:efficientnet-train}, breaklines=true, basicstyle=\footnotesize\ttfamily]
# Estratto da: 04_efficientnet_train_CNN.py
from efficientnet_pytorch import EfficientNet
from torchvision import datasets, transforms
from torch.utils.data import DataLoader

data_dir = r"D:/Thesis/AI Tool Evaluation/efficientnet_dataset"
model_path = "efficientnet_baseline.pth"
batch_size = 32
num_epochs = 10

transform = transforms.Compose([
    transforms.Resize((224, 224)),
    transforms.ToTensor(),
    transforms.Normalize(mean=[0.485, 0.456, 0.406],
                         std=[0.229, 0.224, 0.225]),
])

train_data = datasets.ImageFolder(os.path.join(data_dir, "train"), transform=transform)
val_data = datasets.ImageFolder(os.path.join(data_dir, "val"), transform=transform)

model = EfficientNet.from_name('efficientnet-b0')
model.load_state_dict(torch.load("D:/Thesis/Pretrained/efficientnet-b0-355c32eb.pth"))
model._fc = nn.Linear(model._fc.in_features, len(train_data.classes))
...
# Allenamento, validazione e salvataggio finale
\end{lstlisting}

\subsection*{Inferenza e rilevamento OOD}
\begin{lstlisting}[language=Python, caption={Script per inferenza e rilevamento immagini fuori distribuzione}, label={lst:efficientnet-test}, breaklines=true, basicstyle=\footnotesize\ttfamily]
# Estratto da: 04_efficientnet_ood_test_CNN.py
from efficientnet_pytorch import EfficientNet

model = EfficientNet.from_name('efficientnet-b0')
model._fc = nn.Linear(model._fc.in_features, num_classes)
model.load_state_dict(torch.load(model_path, map_location=device))

# Funzione per determinare se una classe predetta rappresenta un'arma
def is_weapon_class(class_name):
    keywords = ["ak", "glock", "colt", "uzi", "m16", "rifle", "pistol", "gun", "firearm", "revolver", "weapon"]
    return "in-distribution" if any(k in class_name.lower() for k in keywords) else "OOD"
\end{lstlisting}
